\section{A structural transition for shared variables}\label{ct:section}

The trials with large horizons in Lemma~\ref{ctr:reservoir} can use variables
shared across sites.  Our immediate objective is to reorganize one layer
of that shared randomness while preserving the normalized trial bound at
a fixed finite horizon.  The records of conditional laws below quantify the
information needed for the later replacement by independent site
sampling.

\subsection{Continuation laws retained at selected levels}

Lemma~\ref{cs:sensitivity} controls the cost of moving nearby averaging
exponents. Once the exponents on a block have been made equal, its
transforms combine into one conditional logarithmic average. We seek a
probability measure at the block's start from which each fresh site's
continuation can be sampled, conditionally independently of the other
sites. The measures may still depend on the common ancestor data.
Section~\ref{cc:section} will choose finite test families fine enough to
control the error of this replacement.

Fix an exponent width $\eta\in(0,1)$. In an \emph{active} hierarchy,
choose boundary indices
\[
 1=b_0<\cdots<b_v=d.
\]
The first level is reserved for the next reduction, and the consecutive
blocks $(b_{i-1},b_i]$, $1\le i\le v$, are called processed. The case
$v=0$ is allowed when $d=1$. In a \emph{fully processed} hierarchy,
the boundaries instead begin at $b_0=0$, so these blocks cover every
level from $1$ to $d$. Each processed block $(b,c]$ must satisfy
\begin{equation}
 t_c-\eta\le t_\ell\le t_c\qquad(b<\ell\le c).
 \label{ctr:band}
\end{equation}

For a concrete example, consider the last processed block $(b,d]$.
Fix a complete global path $z=(z_0,\ldots,z_d)$ and a single site's
history $u_{0:b}=(u_0,\ldots,u_b)$. Continuing only that site's uniform
sampling gives the message-pair law
\[
 \nu_{z,u_{0:b}}
   =\operatorname{Law}_{u_{b+1:d}}(f^+,f^-)
       \quad\hbox{on }[0,1]^2.
\]
At boundary $b$ we store a probability measure $S_b$ that uses only the
global and site history through $b$. We ask it to predict finitely many
continuous test means of $\nu_{z,u_{0:b}}$, with a small mean-square error
over $u_{0:b}$ for almost every fixed $z$. Thus the estimate averages
ancestor site variables but does not average away dependence on later
global variables. Earlier blocks require the same construction for the
law of a probability-valued object, which leads to the iterated spaces
below.

Let $K^{(0)}=[0,1]^2$ with the sup metric and, recursively, let
\[
 K^{(j+1)}=\mathcal P(K^{(j)})
\]
with the Wasserstein--$1$ metric induced by the preceding metric. These
spaces are compact and have diameter at most one. At the bottom boundary
put $S_d=(f^+,f^-)\in K^{(0)}$. Working upward through the boundaries,
if $S_c$ has values in $K^{(j)}$, assign to $S_b$ values in $K^{(j+1)}$.
Each $S_b$ is a measurable function only of the globals and the single
site's generators through level $b$. A type is needed only at a boundary.
The assigned type need not equal a conditional law exactly.

Choose finite nonempty families $\Phi_j$ of continuous functions
$K^{(j)}\to[0,1]$ and a tolerance $\delta>0$. For a processed block
$(b,c]$ with $S_c\in K^{(j)}$, the required record is
\begin{equation}
 \int\left|
     \int\phi(S_c)\,\dd u_{b+1:c}-S_b(\phi)
     \right|^2\dd u_{0:b}\le\delta^2
       \qquad(\phi\in\Phi_j),
 \label{ctr:record}
\end{equation}
for almost every fixed complete global path, including its root.
Here $u_{p:q}$ lists one site's generators $u_p,\ldots,u_q$, all site
integrals are Lebesgue integrals, and $S_b(\phi)$ is the integral of
$\phi$ against the probability measure $S_b$.

The pathwise qualification preserves this record under the posterior of
Lemma~\ref{ctr:posterior} and its additional positive first-level tilt.
We keep the type functions unchanged, ignoring newly added global
coordinates. Fresh site generators retain their product uniform law,
and the enlarged global law is absolutely continuous with respect to
the old law. The original exceptional set of complete global paths
therefore remains null; the types are not recomputed as conditional
expectations under the new posterior.

The depth-one reservoir trial starts active with boundary $1=d$ and no
processed blocks. The next transition adds one processed block while
preserving every previous record. Its bound on inserted levels and its
decrease of the first remaining exponent will force termination after
finitely many reductions.

\subsection{One step of the reduction}

The following result turns a sequence of trials with growing horizons
into a trial at a fixed horizon, with one more recorded block.  Its input depth, boundaries,
test functions, and output horizon are fixed before the input horizon
tends to infinity.

\begin{lemma}[Structural transition]
\label{ct:transition}
Fix a density $a>0$, a finite penalty $\beta>0$, $\eta\in(0,1)$,
$\delta>0$, and $m_0>0$. Suppose there is a
sequence of active trials of fixed depth $d$ and fixed boundary indices,
with horizons $N\to\infty$, satisfying
\begin{equation}
 \frac{G_N}{N}\le U+o(1),\qquad t_1\ge m_0,
 \label{ct:input-bound}
\end{equation}
for a finite $U$. Suppose the old blocks satisfy
Equations~\eqref{ctr:band} and~\eqref{ctr:record} with fixed finite
continuous test families. Fix also a finite nonempty continuous
$[0,1]$-valued test family on the type space of the current $S_1$.
After passing to a subsequence along which all old exponents
converge, for each fixed positive integer $k$ there is a finite-depth
trial of horizon $k$ with $G_k/k\le U$ having the following
properties.

For some $h\in\{0,1,2\}$, the output has depth $d+h$: the $h$ new
levels are inserted above the old first level, which becomes the
\emph{sample level} $h+1$.  Every old block $(b,c]$ becomes
$(b+h,c+h]$ and retains its exponent band, test family, tolerance, and record.
One new processed block ends at level $h+1$ and satisfies
Equations~\eqref{ctr:band} and~\eqref{ctr:record} for the specified
new test family. Either this new block starts at the root, making
the hierarchy fully processed, or it starts at the new first
level and the hierarchy remains active, with first exponent
\begin{equation}
 0<m_{\mathrm{new}}\le(1-\eta)\lim t_1.
 \label{ct:progress}
\end{equation}
All records hold for almost every complete global path.
\end{lemma}

The subsequent conditional-law and array subsequences may depend on the
fixed output horizon $k$.  The conclusion asserts no uniform rate as $k$
varies; this is the order of limits used in the finite iteration later.


\paragraph{Proof strategy.}
In Subsections~\ref{ct:sub-probing}--\ref{ct:sub-selection}, we perturb
the law at the initial level by a compensated Gaussian field.  After selecting
one conditional law, this perturbation
gives identities for inner products of conditional site-test means.
The identities remain valid when the test also observes all retained
descendant and site data.  Subsection~\ref{ct:sub-cuts} uses these identities
to organize the samples into a nested
partition at at most two thresholds.  An exchangeability theorem then
represents the partitioned array using independent variables along a
sample tree and a site tree, including the variables shared by all sites.

Subsection~\ref{ct:sub-endpoint} identifies the endpoint transforms by
replacing each conditional
integral temporarily with a finite average over independent children.
This also determines the new exponents.  Finally, in
Subsection~\ref{ct:sub-records}, empirical averages over
fresh sites transfer closed inner-product constraints.  A Hilbert-space
support argument converts them into records for almost every complete
shared path.

Throughout the proof the old depth, boundary indices, finite test
families, and desired output horizon $k$ are fixed.  Write $m=t_1$ for
the first input exponent.  Along the input sequence $N\to\infty$,
$m\ge m_0$, and Equation~\eqref{ct:input-bound} holds.  Pass to a
subsequence on which all old exponents converge. We suppress the sequence
index: before the weak limit below, $m$ and $t_j$ denote the input
exponents; afterward they denote their limits. In both uses $m\ge m_0$.
Expectations over the Gaussian fields below are denoted by $\E_g$.
Their base probability measure is fixed before that expectation is
taken.

\subsection{Probing the first conditional site law}\label{ct:sub-probing}

We first express the effect of the first shared variable on the chosen
site tests.  Fix the root data for the moment.
Let $\Phi=\{\phi_1,\ldots,\phi_{J_f}\}$, with $J_f\ge1$, be the finite test family
for the current first-boundary type $S_1$.  Root dependence is suppressed
in the notation.  Define
\begin{equation}\label{ct:features}
 H_j(z_1,u_0)=\int\phi_j(S_1)\,\dd u_1,
 \qquad
 Q_j(z_1,z'_1)=\int H_j(z_1,u_0)H_j(z'_1,u_0)\,\dd u_0.
\end{equation}
These are Gram kernels of $[0,1]$-valued functions in the separable
space $L^2([0,1])$.  Enumerate as $(C_l)$ all nonconstant monomials in
$Q_1,\ldots,Q_{J_f}$, with nonnegative integer powers.  Tensor products of
their feature vectors show that these too are positive semidefinite,
with diagonals at most one.

Each $Q_j$ is an inner product of two conditional test means.  We
will derive identities for all these overlaps jointly, allowing the test
also to observe the retained descendant data.  This stronger test class
will later be needed when we rank clusters by their masses.  The next
lemma supplies the identities uniformly over the conditional law of the
first-level index.

The identities belong to the Ghirlanda--Guerra family of conditional
overlap identities~\cite[arXiv version, Equation~(28)]{GhirlandaGuerra1998}.
The use of tests that also observe the retained descendant and site data
follows the spin-tested
formulation of Panchenko~\cite[Equation~(14)]{Panchenko2015HE}.
We prove the conditional error estimate needed for the selected laws below.

\begin{lemma}[Conditional Gaussian probing]\label{ct:gaussian}
Let $\mu$ be a probability measure on a standard Borel space.  Let
$(C_l)_{l\ge1}$ be positive semidefinite measurable kernels with separable
Hilbert-space feature representations and $0\le C_l(\zeta,\zeta)\le1$.
Put $a_l=2^{-l}$ and let $g_l$ be independent centered Gaussian fields with
covariances $C_l$.  For $s\ge1$ and $x_l\in[1,2]$, set
\begin{equation}\label{ct:gaussian-tilt}
 h(\zeta)=s\sum_{l\ge1}a_lx_lg_l(\zeta)
       -\frac{s^2}{2}\sum_{l\ge1}a_l^2x_l^2 C_l(\zeta,\zeta),
 \qquad
 \mu^h(\dd\zeta)=\frac{e^{h(\zeta)}\mu(\dd\zeta)}{\mu(e^h)}.
\end{equation}
The normalizer is finite and positive almost surely.
Brackets denote conditionally independent samples from $\mu^h$ given the
Gaussian fields, with any additional marks sampled by a kernel depending
on those samples and with no further dependence on the fields.  Define
\[
 V_l(\zeta)=\frac{g_l(\zeta)}{s a_lx_l}-C_l(\zeta,\zeta),
 \qquad
 e_l=\E_g\left\langle
  \left|V_l(\zeta^1)-\E_g\langle V_l(\zeta^1)\rangle\right|
 \right\rangle.
\]
Uniformly in $\mu$, in the kernels, and in the other parameters,
\begin{equation}\label{ct:error-integral}
                         \int_1^2 e_l\,\dd x_l=o_l(1)
                         \qquad(s\longrightarrow\infty).
\end{equation}
For every integer $n\ge1$ and every bounded measurable $f$, $|f|\le1$, of the first $n$ samples and
their marks,
\begin{align}\label{ct:approx-gg}
 \bigg|n\E_g\langle fC_l(1,n+1)\rangle
   -\sum_{i=2}^n\E_g\langle fC_l(1,i)\rangle
   -\E_g\langle f\rangle\E_g\langle C_l(1,2)\rangle\bigg|
 \le e_l.
\end{align}
Here $C_l(i,j)=C_l(\zeta^i,\zeta^j)$.
\end{lemma}

\begin{proof}
Choose measurable feature maps $v_j$ with
$C_j(\zeta,\zeta')=\langle v_j(\zeta),v_j(\zeta')\rangle$; their
norms are at most one.  Fix $l$ and write $x=x_l$.  We allow $x$ to range over
$[0,3]$ during the proof.  Numerical constants denoted by $K_l$
depend only on $l$.  Consider
\[
 \Psi(x)=\frac{1}{s^2a_l^2}\log\mu(e^h)+\frac{x^2}{2},
 \qquad
 A_x(\zeta)=\frac{g_l(\zeta)}{sa_l}-xC_l(\zeta,\zeta).
\]
Differentiation gives
\begin{equation}\label{ct:convex-derivatives}
 \Psi'(x)=\langle A_x\rangle+x,
 \qquad
 \Psi''(x)=s^2a_l^2\langle(A_x-\langle A_x\rangle)^2\rangle
                 +1-\langle C_l(\zeta,\zeta)\rangle\ge0.
\end{equation}
The deterministic diagonal correction in \eqref{ct:gaussian-tilt}
is exactly what permits both this convexity and the cancellation in
\eqref{ct:approx-gg} below.

For every fixed $\zeta$, $\E_g e^{h(\zeta)}=1$ and
$-Cs^2\le\E_g h(\zeta)\le0$, with an absolute constant $C$.
The two Jensen inequalities therefore give
\begin{equation}\label{ct:log-normalizer-bounds}
       -Cs^2\le\E_g\log\mu(e^h)\le0.
\end{equation}
Consequently the convex function $\bar\Psi=\E_g\Psi$ is bounded on
$[0,3]$ by a constant depending only on $l$, and its one-sided
derivatives are uniformly bounded on $[1/2,5/2]$.  Integrating
\eqref{ct:convex-derivatives} over $[1,2]$ yields
\begin{equation}\label{ct:thermal-variance}
 \int_1^2\E_g\langle(A_x-\langle A_x\rangle)^2\rangle\,\dd x
                         \le K_l s^{-2}.
\end{equation}
Interchanging expectation with the first derivative is legitimate here:
on any compact interior interval, convexity bounds the derivative by
secants to two fixed outer points, whose absolute values are integrable
by \eqref{ct:log-normalizer-bounds} and Gaussian integrability.

The Gaussian gradient of $\log\mu(e^h)$ has norm at most $Cs$.
Indeed it is the Gibbs mean of the direct-sum feature vector
$(s a_jx_j v_j(\zeta))_j$, whose squared norm is at most
$s^2\sum_j a_j^2x_j^2$.  The Gaussian variance inequality gives
\begin{equation}\label{ct:function-concentration}
                  \E_g|\Psi(x)-\bar\Psi(x)|\le K_l/s.
\end{equation}
For $0<u<1/2$, compare $\Psi'(x)$ with its left and right secants of
length $u$, and do the same for $\bar\Psi$.  After integrating over
$x\in[1,2]$, \eqref{ct:function-concentration} bounds the random
secant errors by $K_l/(su)$.  The remaining deterministic error is at
most
\[
 \int_1^2\bigl(\bar\Psi'_+(x+u)-\bar\Psi'_-(x-u)\bigr)\,\dd x
 \le K_lu;
\]
the integral telescopes and the interior derivatives are bounded.
Thus
\begin{equation}\label{ct:derivative-concentration}
 \int_1^2\E_g|\Psi'(x)-\E_g\Psi'(x)|\,\dd x
                         \le K_l\bigl((su)^{-1}+u\bigr).
\end{equation}
Since $V_l=A_x/x$ and $x\ge1$, the sum of
\eqref{ct:thermal-variance}, by Cauchy--Schwarz, and
\eqref{ct:derivative-concentration}, with $u=s^{-1/2}$, proves
\eqref{ct:error-integral}.

For completeness, all the preceding operations can first be performed
with finitely many Gaussian coordinates in finitely many feature spaces.
The resulting bounds do not depend on these truncations.  The Gaussian
series converge in $L^2(\mathbb P_g\otimes\mu)$.  For fixed $s$ their
limit is finite for $\mu$-almost every index, for almost every field
realization, so the normalizer is positive.  Also
$\E_g\mu(e^h)=1$ by Fubini, so the normalizer is finite almost surely.
For fixed $s$ the
exponentials have uniformly bounded moments of every fixed positive
order, while Jensen's inequality gives
\[
      \mu(e^h)^{-q}\le\mu(e^{-qh})\qquad(q>0).
\]
The right side has uniformly bounded Gaussian expectation.  These
bounds give uniform integrability of the logarithms and of the replica
integrals, including polynomial factors in $g_l$.  They justify passage
from finite-dimensional differentiation, concentration, and Gaussian
integration by parts to the displayed formulas.  Equivalently one can
construct the fields by isonormal processes on the separable feature
spaces; joint measurability is needed only outside a
$\mathbb P_g\otimes\mu$ null set.

Average the additional marks conditional on their sample indices.  The
resulting test is independent of the fields, so Gaussian integration by
parts gives
\[
 \E_g\left\langle f\frac{g_l(\zeta^1)}{s a_lx_l}\right\rangle
  =\E_g\left\langle
       f\left(\sum_{i=1}^n C_l(1,i)-nC_l(1,n+1)\right)
    \right\rangle.
\]
Subtracting the diagonal term in $V_l$ removes $C_l(1,1)$.
The case $f=1$, $n=1$ says
$\E_g\langle V_l\rangle=-\E_g\langle C_l(1,2)\rangle$.
Finally,
\[
 \left|\E_g\langle fV_l\rangle
       -\E_g\langle f\rangle\E_g\langle V_l\rangle\right|
 \le e_l,
\]
which proves \eqref{ct:approx-gg}.
\end{proof}

\subsection{Selecting one law and retaining its site data}\label{ct:sub-selection}

For our concrete kernels the feature spaces are tensor powers of
$L^2([0,1])$.  Fixed orthonormal bases in these spaces construct the
Gaussian fields measurably also when the root global is random, before
we condition on its value.

Apply Lemma~\ref{ct:gaussian} with $s=N^{1/4}$.  Add $h(z_1)/m$
to both endpoint losses in the horizon-$N$ comparison.  This term
passes unchanged through the transforms below the first level.
The fixed-$s$ exponential and inverse-normalizer bounds in the proof of
Lemma~\ref{ct:gaussian} make the normalizer finite and positive almost
surely and its logarithm integrable.  Together with the conditional packet
bound, these facts verify the integrability requirements for the added
term in Lemma~\ref{ctr:posterior}.
Conditioning on all other endpoint data, its change to the first
transform is
\[
                     \frac1m\log\mu_*(e^h),
\]
where $\mu_*$ is the unperturbed first-level posterior.  By
\eqref{ct:log-normalizer-bounds}, the absolute change of its
Gaussian-averaged value is at
most $Cs^2/m$.  The total perturbation error is therefore $o(N)$,
uniformly while $m\ge m_0$.

Telescope the comparison in complete blocks of $k$ packets, with a
remainder of fewer than $k$ packets, using
Lemma~\ref{ctr:posterior}.  For a block, condition on all the
base data belonging to the other packets, including their root-site
variables and the original root global.  Before the Gaussian tilt,
its enlarged first-level law is then a fixed probability measure
$\mu$ on an index $\zeta$.  This index includes the other packets'
first-level variables.  The kernels in \eqref{ct:features} use its
original $z_1$ projection only: the fresh site's variables remain
independently uniform.  All later posterior kernels depend on these
base data and on their histories, but not otherwise on the Gaussian
fields.  Thus the marks used below have exactly the independence
required in Lemma~\ref{ct:gaussian}.

Choose a complete block uniformly.  Its posterior $k$-packet
functional, still averaged over Gaussian fields and fresh data,
satisfies
\begin{equation}\label{ct:selected-value}
                      \E_g G_k^{\rm post}\le kU+o(1)
\end{equation}
on average over the block and base data and over independent uniform
$x_l\in[1,2]$.  This follows from \eqref{ct:input-bound}, the
$o(N)$ endpoint error, and the $O_{a,\beta}(k)$ remainder bound.
Furthermore, conditionally on any such base data and parameters,
\begin{equation}\label{ct:conditional-value-bound}
                    |\E_g G_k^{\rm post}|\le C_{a,\beta}k.
\end{equation}
Indeed a normalized logarithmic transform lies between the extrema
of its argument, and each fresh packet has the root bound
$|X|\le\log2+\beta D$.

The uniform bounds \eqref{ct:error-integral}, Markov's inequality,
and a diagonal choice yield sets of choices of block, base data, and parameters of
probability tending to one on which $e_l\to0$ for every fixed $l$.
More explicitly, choose a number of initial kernels tending to infinity
so slowly that their summed mean errors tend to zero, and require that
their sum of errors be at most the square root of that mean.  Discard
also the conditioning null sets on which a structural record fails.
By \eqref{ct:conditional-value-bound}, removal of the discarded
set changes the mean in \eqref{ct:selected-value} by $o(1)$.
For each $N$ we may consequently choose deterministic block, base data,
and parameters so that both \eqref{ct:selected-value} and all these
vanishing-error conditions hold.  The Gaussian fields remain averaged.

Under each selected conditional law, keep the Gaussian fields random.
Conditional on them, sample countably many first-level indices independently
from $\mu^h$, in their natural order $u=1,2,\ldots$. A natural label path
has the form
\[
 (u,\xi),\qquad u\in\N,\quad\xi\in\N^{d-1},
\]
where $\xi$ indexes the old descendant hierarchy below the first-level
sample; for $d=1$, it is the empty label. Draw independent fresh root-site uniforms for the sites $i\in\N$, with
each site's root uniform shared by all its label paths. Conditional on
the selected base and Gaussian fields, draw first-level site uniforms
independently across sites and first-level sample indices. Below each sample,
continue the old descendant hierarchy using its conditional global kernels
and independent site uniforms. At each branch, distinct children use independent generators conditional
on the parent history. Each site's root uniform is reused on all label
paths; its later site uniforms are reused along their common label
prefixes.

Let $\mathcal B$ be the old boundary set, including $1$ and $d$. Retain
at each path and site the decoration
\[
 D_{u,\xi,i}=(S_b^{u,\xi,i})_{b\in\mathcal B},
\]
repeating each boundary value on its descendants. Also retain all pair
kernels $(Q_j(u,v))$. Let $\mathcal L_N$ be the joint law of these
countable arrays. This law averages the Gaussian fields and all the newly
sampled data; the selected block, old-packet base data, and perturbation
parameters are fixed.

Each decoration coordinate lies in one of the compact type spaces used
in Equation~\eqref{ctr:record}, and every kernel coordinate lies in
$[0,1]$. The countable product is compact and metrizable, so along a
subsequence $\mathcal L_N$ converges weakly to a law $\mathcal L$.
The product term in \eqref{ct:approx-gg} is a product of expectations
under $\mathcal L_N$ and therefore converges to the corresponding product
under this one limiting law. Write $\E$ for expectation under
$\mathcal L$ in the identities below.

For $n\ge1$, let $\mathcal A_n$ be the sigma-field generated by
$D_{u,\xi,i}$ for $u\le n$ and every $\xi,i$, together with
$Q_j(u,v)$ for $u,v\le n$ and $1\le j\le J_f$. For every bounded $\mathcal A_n$-measurable $f$ and every
bounded Borel function $\psi$ on $[0,1]^{J_f}$, the limit satisfies
\begin{equation}\label{ct:marked-gg}
 n\E[f\psi(\mathbf Q(1,n+1))]
 =\E[f]\E[\psi(\mathbf Q(1,2))]
       +\sum_{i=2}^n\E[f\psi(\mathbf Q(1,i))],
 \qquad \mathbf Q=(Q_1,\ldots,Q_{J_f}).
\end{equation}
Indeed, first take $f$ to be a continuous cylinder function of those
decorations and kernels. It is an allowable mark-tested function in
Lemma~\ref{ct:gaussian}. Weak convergence gives the identity for
nonconstant monomials $\psi$; constants satisfy it directly. Polynomial
approximation gives it for continuous $\psi$. Monotone-class arguments
extend it first to bounded Borel $\psi$ and then to every bounded
$\mathcal A_n$-measurable $f$.
The joint law of the decorations and pair kernels is invariant under
simultaneous permutations of the first-level sample indices. The
hierarchical permutation invariances of the old descendant trees and
the simultaneous permutation invariance of the sites also pass to the
limit.

\subsection{Ordered overlaps and marked finite cuts}\label{ct:sub-cuts}

We next turn the identities into a nested partition of the samples.
The external ultrametricity theorem is stated for actual inner products
of Hilbert-space samples.  The following observation records the passage
from our arrays to that setting, including the diagonal.

\begin{samepage}
\begin{lemma}[Ultrametricity from off-diagonal identities]\label{ct:gram-bridge}
Let $R=(R_{ij})_{i,j\ge1}$ be a symmetric, weakly exchangeable positive
semidefinite array with $R_{ii}=1$.  Suppose that for every $n\ge2$,
every bounded Borel function $f$ of the off-diagonal entries among the
first $n$ indices, and every bounded Borel function $\psi$, one has
\[
 n\E[f\psi(R_{1,n+1})]
 =\E[f]\E[\psi(R_{12})]+\sum_{i=2}^n\E[f\psi(R_{1i})].
\]
Then $R_{12}\ge\min\{R_{13},R_{23}\}$ almost surely, and the analogous
inequality holds for every triple of distinct indices.
\end{lemma}
\end{samepage}

\begin{proof}
The Dovbysh--Sudakov representation, in Proposition~1 of
\cite{Panchenko2010DS}, realizes the array as a mixture of arrays
\[
 R_{ij}=\langle v_i,v_j\rangle+a_i\1_{\{i=j\}},\qquad a_i\ge0,
\]
where the pairs $(v_i,a_i)$ are conditionally independent with the
same law on a separable Hilbert space times $[0,\infty)$.  The unit
diagonal puts every $v_i$ in the unit ball.  Let $\mathcal G$ denote
the random directing law of the vectors and put
$q_* = \max\supp\operatorname{Law}(R_{12})$; this support is compact
because positive semidefiniteness and the unit diagonal imply
$|R_{ij}|\le1$.

Theorem~2(a) of Panchenko~\cite{Panchenko2010GG} applies to these
off-diagonal Ghirlanda--Guerra identities and gives
\[
 \mathcal G\{v:\|v\|^2=q_*\}=1\qquad\hbox{almost surely}.
\]
Its sphere conclusion has no finite-support hypothesis.  Now form the
true Gram array $\widehat R_{ij}=\langle v_i,v_j\rangle$, including its
diagonal.  It has the same off-diagonal entries as $R$ and the
deterministic diagonal $q_*$.  A bounded test of its complete first-$n$
matrix is therefore a bounded test of the off-diagonal matrix with
constants substituted on the diagonal.  Thus the displayed identities
give the full Ghirlanda--Guerra identities for $\widehat R$ under the
joint, averaged law.  No conditioning on a realization of
$\mathcal G$ is used in this implication.

The random measure $\mathcal G$ is supported on the unit ball of a
separable Hilbert space.  Theorem~1 of
Panchenko~\cite{Panchenko2013UM} now applies to its true Gram array and
gives the asserted ultrametric inequalities for the off-diagonal
entries.  This theorem likewise permits a general overlap distribution.
\end{proof}

\begin{lemma}[Ordering of the vector overlaps]\label{ct:ordering}
The support $\mathcal S$ of the off-diagonal law of $\mathbf Q$ in
\eqref{ct:marked-gg} is totally ordered componentwise.  The scalar off-diagonal overlaps
\[
 Q(u,v)=J_f^{-1}\sum_{j=1}^{J_f}Q_j(u,v),\qquad u\ne v,
\]
satisfy the ultrametric inequalities on triples of distinct indices, and,
for $v,v'\in\mathcal S$,
\begin{equation}\label{ct:coordinate-control}
             |v_j-v'_j|\le J_f\left|J_f^{-1}\sum_i(v_i-v'_i)\right|.
\end{equation}
\end{lemma}

\begin{proof}
In each replica array, set $Q_j(u,u)=1$ for every sample index $u$.
Leave all $Q_j(u,v)$ with $u\ne v$ unchanged, even when the sampled
values of $z_1$ coincide.  Before passage to the limit this adds a
nonnegative diagonal matrix to a Gram matrix, so it
preserves positive semidefiniteness.  The finite-matrix condition is
closed and hence also holds in the limit.  Every convex combination
with positive rational coefficients is therefore a weakly exchangeable
positive semidefinite array with unit diagonal.  Its off-diagonal
identities follow from \eqref{ct:marked-gg}, so
Lemma~\ref{ct:gram-bridge} makes it ultrametric.

There are countably many rational combinations, so their conclusions
hold simultaneously.  In any off-diagonal triple the three vector
overlaps must contain an equal pair.  Otherwise some positive
rational linear functional avoids their three equality hyperplanes
and gives three distinct scalar overlaps, contrary to ultrametricity.
If the repeated vector is $v$ and the remaining vector is $v'$, then
$v\le v'$ componentwise: otherwise a positive rational functional
concentrated on a violating coordinate would make the unique value
the strict minimum, again impossible.

It follows that $\mathbf Q(1,2)$ and $\mathbf Q(1,3)$ are almost
surely comparable.  If two points in $\mathcal S$ were incomparable,
they would have incomparable neighborhoods $A,B$ of positive
pair-law probability.  Apply \eqref{ct:marked-gg} with $n=2$,
$f=\1_{\{\mathbf Q(1,2)\in A\}}$, and
$\psi=\1_B$.  Since $A\cap B=\varnothing$, it gives probability
$\Prob(\mathbf Q\in A)\Prob(\mathbf Q\in B)/2>0$ of seeing an
incomparable pair, a contradiction.  Componentwise ordering implies
\eqref{ct:coordinate-control}; ultrametricity of $Q$ was already
obtained from the equal-weight combination.
\end{proof}

For the rest of this subsection, $Q$ denotes a generic off-diagonal
overlap, with the law of $Q(1,2)$.  Fix the exponent width $\eta\in(0,1)$ and a number $\gamma>0$ to be
chosen below.  Let
\[
 q=\min\{t:\Prob(Q\le t)\ge1-\eta\},
 \qquad x_p=\Prob(Q<q)\le1-\eta.
\]
We choose $r>q$ so that
\begin{equation}\label{ct:quantile-band}
 x_c=\Prob(Q<r)>x_p,\qquad x_c\ge1-\eta,
 \qquad \operatorname{diam}\supp(Q\mid q\le Q<r)\le\gamma.
\end{equation}
Here the conditioning event has positive probability.  If
$\Prob(q\le Q<q+\gamma)>0$, take $r=q+\gamma$.  Otherwise $q$ is
not an atom, $\Prob(Q\le q)=1-\eta$, and there is a gap of length
at least $\gamma$ to its right.  Let $q'$ be the least support point
of the remaining mass in $[q+\gamma,\infty)$ and take
$r=q'+\gamma$.  Then the band $[q,r)$ has its actual support in
$[q',q'+\gamma]$, and \eqref{ct:quantile-band} follows.  This
also covers endpoint atoms: $r$ may exceed the maximum overlap.

Insert the parent relation $Q(u,v)\ge q$ if $x_p>0$, and the child
relation $Q(u,v)\ge r$ if $x_c<1$.  Make each relation reflexive on
the diagonal.  Ultrametricity makes them nested equivalence
relations.  The proper cuts have parameters
\begin{equation}\label{ct:cut-parameters}
                       0<x_1<\cdots<x_h<1,
                       \qquad 0\le h\le2.
\end{equation}
All inequalities at a cut are non-strict, so atoms are included
without ambiguity.
The partition records only membership in these two relations.  The
overlap law within a band may still have infinitely many support points;
the later record estimate will use the support diameter in
Equation~\eqref{ct:quantile-band}.
Figure~\ref{fig:hidden-cuts} illustrates the two-cut case.
\begin{figure}[!htbp]
\centering
\begin{tikzpicture}[>=Latex,x=1cm,y=1cm,every node/.style={font=\small}]
 \draw[->] (0,0.45)--(0,-4.45);
 \node[anchor=east] at (-.2,.35) {root};
 \foreach \y/\label in {-.65/{parent cut},-1.85/{child cut},-3.05/{old sample level},-4.1/{old descendants}} {
  \draw (-.08,\y)--(.08,\y);
  \node[anchor=east] at (-.2,\y) {\label};
 }
 \node[anchor=west] at (.25,-.65) {$m x_p\le(1-\eta)m$};
 \node[anchor=west] at (.25,-1.85) {$m x_c\ge(1-\eta)m$};
 \node[anchor=west] at (.25,-3.05) {$m$};
 \node[anchor=west] at (.25,-4.1) {$t_2,\ldots,t_d$};
 \draw[decorate,decoration={brace,amplitude=5pt}] (4.3,-1.65)--(4.3,-3.25);
 \node[anchor=west,align=left,text width=3.1cm] at (4.6,-2.45)
  {new processed block:\\exponents lie in $[m-\eta,m]$};
\end{tikzpicture}
\caption{One structural transition when both cuts are proper. The parent
exponent decreases by a fixed factor; the levels strictly below it through
the old sample level form the new processed block. The cut parameters are
$x_p=\Prob(Q<q)$ and $x_c=\Prob(Q<r)$. Trivial parent or child cuts are
omitted, as specified in the proof. The old descendant row is absent when
$d=1$.}
\label{fig:hidden-cuts}
\end{figure}


Panchenko's pure-state theorem proves hierarchical exchangeability of spin
arrays and their independence from cluster weights under spin-tested
Ghirlanda--Guerra identities~\cite[Theorem~1]{Panchenko2015HE}.
To describe the corresponding indexing here, rank cut clusters by decreasing
frequency within their parents, resolving ties by the least natural sample
index, and list samples within each deepest cluster in natural order.
Write $(\alpha,l)\in\N^h\times\N$ for a sample's new address: $\alpha$
is its hidden-cluster address and $l$ its number within that cluster.
The decoration array becomes $D'_{\alpha,l,\xi,i}$. The whole old
descendant subtree moves with its first-level sample $u$; the site index
$i$ stays common. The following local argument proves independence of
this reindexed array from the full partition that determines its addresses.

\FloatBarrier
\begin{lemma}[The marked cut representation]\label{ct:marked-cuts}
The nested cut partition has the sample-partition law of an RPC with
parameters \eqref{ct:cut-parameters}. The reindexed decorations
$D'_{\alpha,l,\xi,i}$ are independent of the full partition and are
hierarchically exchangeable in the hidden-cluster and sample coordinates.
They retain the independent hierarchical exchangeabilities of the old
descendant trees and the simultaneous exchangeability of the sites.
\end{lemma}

\begin{proof}
For a cut of parameter $x_j$, let $C$ be a cluster seen among the
first $n$ samples, of size $n_C$.  Select an anchor in $C$ by a
finite case distinction.  Equation~\eqref{ct:marked-gg}, with
the indicator of the anchor's cut relation as $\psi$, shows that,
conditionally on all past pair and decoration data, the next
sample belongs to $C$ with probability
\begin{equation}\label{ct:partition-prediction}
                            \frac{n_C-x_j}{n}.
\end{equation}
Indeed the independent term contributes $(1-x_j)/n$ and the past
anchor-neighbor terms contribute $(n_C-1)/n$.  Differences between
these probabilities specify the creation of new children within
each parent.  These are exactly the RPC predictions in
Lemma~\ref{ctr:rpc}, Equation~\eqref{ctr:prediction}.  They determine every
finite nested
partition law.  Consequently the full partition has that law, its
clusters have positive frequencies, and the frequencies at every
level have total mass one.  The ranked leaf frequencies have the
corresponding RPC distribution.  For $h=0$ these statements mean a
single cluster of mass one.

We now prove independence from the full partition using the decoration
tests in Equation~\eqref{ct:marked-gg}.  Let $I$ be a deterministic
ordered tuple of distinct sample indices, let $D_I$ be all their
decorations, and let $\Pi$ be the full nested partition.  Expose $I$
first, using exchangeability.  Given its cut pattern and decorations,
iterating \eqref{ct:partition-prediction} shows that the conditional
law of the partition on all remaining indices depends only on that
cut pattern.  In this iteration one can condition down from all
intermediate decorations, since the predictions do not depend on them.
Extending from finite restrictions to the countable partition gives,
for every bounded decoration test $f$,
\begin{equation}\label{ct:conditional-decoration-law}
 \E[f(D_I)\mid\Pi]=K_{\Pi|_I}(f),
\end{equation}
where the probability kernel $K$ depends only on the ordered internal
cut pattern $\Pi|_I$, and not on the actual values of the indices.

The inverse image of a fixed finite tuple of distinct positions in
the ranked-cluster/natural-order indexing is a tuple of original
indices measurable with respect to the partition.  Partitioning
according to its countably many possible values and applying
Equation~\eqref{ct:conditional-decoration-law} proves that the reindexed tuple's decoration
law is independent of the partition.  Its internal cut pattern is
fixed by its new positions.  The same argument shows invariance
under ancestry-preserving permutations of the hidden clusters and
of the sample indices within them.  It applies to arbitrary finite
decoration cylinders and hence to the complete arrays.  The old
descendant permutations and the site permutations do not change
the cut partition, so their invariances are retained as well.
\end{proof}

We have now separated the RPC partition from the complete decoration
array.  The next representation supplies path variables for that array.

Combine the $h$ hidden levels, the sample level, and the $d-1$ old
descendant levels into one label tree of depth $h+d$.  Together with
the depth-one site tree, the decorations form an array on a product
of two finite-depth, countably branching trees.  Their entry space,
a finite product of compact type spaces, is standard Borel.
The invariances just proved give the separate ancestry-preserving
permutation invariances for these two trees.
Theorem~2 of Austin and Panchenko~\cite{AustinPanchenko2014} therefore
represents this array by measurable functions of independent uniform
variables indexed along the product of the two ancestor paths.
Those paired with the site-tree root are shared globals $z_l$;
those paired with site $i$ are site variables $u_{l,i}$.
We take this representation independent of the partition, as
permitted by Lemma~\ref{ct:marked-cuts}.  All shared globals are
retained, including the root global.  The representation concerns the
site array alone; fresh packet counts and signs are still drawn
independently at the root, as in Section~\ref{ctr:section}.

An old type at a boundary shifted by $h$ has a version using only
the variables through that boundary.  To see this, take two leaves
agreeing through the boundary and diverging immediately afterward.
Their boundary-type entries, for the same site, are equal almost
surely by the duplicated-entry property.  In the representation
these two entries are conditionally independent and identically
distributed given the prefix variables.  Their conditional law is
therefore a point mass.  A measurable version depends only on the
prefix.  There are only countably many entries, so these versions
agree simultaneously throughout the array.  At the last boundary
there is no reduction to make.

\subsection{Identifying the endpoint functional}\label{ct:sub-endpoint}

The representation just obtained has new levels.  We determine their
exponents and their endpoint values from the original natural-order
sampling.  This calculation uses the full array law behind the finite
partition, including the distribution within its bands.

\begin{lemma}[Uniform empirical tree evaluation]\label{ct:empirical}
Let $C\ge0$, let $B\ge1$ be an integer, and let $|X|\le C$ in a hierarchy of fixed depth $d$ with
exponents in $(0,1]$.  At every vertex replace its conditional integral
by the empirical average over $B$ conditionally independent children.
The resulting nested logarithmic transform differs from the exact
transform in mean absolute value by at most
\[
                            A(C,d)B^{-1/2},
\]
uniformly in the conditional sampling kernels, the exponents, and
the root data.  The same statement holds for finitely many sites
sampled jointly at every vertex.
\end{lemma}

\begin{proof}
For a one-step transform of a random variable $Y\in[-C,C]$,
\[
 T_tY=\frac1t\log\E e^{tY},
 \qquad
 \widehat T_tY=\frac1t\log\left(\frac1B\sum_{b=1}^B e^{tY_b}\right).
\]
The standard deviation of $e^{tY}$ is at most $2tCe^C$, while
$t^{-1}\log v$ is $e^C/t$-Lipschitz on
$[e^{-tC},e^{tC}]$.  Hence
\begin{equation}\label{ct:one-step-empirical}
                      \E|\widehat T_tY-T_tY|
                              \le2Ce^{2C}/\sqrt B.
\end{equation}
If exact child values $Y$ are replaced by values $Y'$ in $[-C,C]$,
interpolation between the two sets of values gives
\[
               |T_tY-T_tY'|\le e^{2C}\E|Y-Y'|,
\]
and the identical bound with an empirical mean holds for an
empirical transform.  Apply these bounds recursively from the
leaves upward.  Every intermediate value remains in $[-C,C]$,
so the constants depend only on $C,d$.  A joint child variable
may include the global coordinate and all the finitely many site
coordinates; conditional independence between children is all that
was used.
\end{proof}

\begin{lemma}[Endpoint passage]\label{ct:endpoint}
Equip the represented hierarchy with exponents
\begin{equation}\label{ct:new-exponents}
                  mx_1,\ldots,mx_h,\ m,\ t_2,\ldots,t_d.
\end{equation}
For each endpoint of horizon $k$, its expected nested transform is
the subsequential limit of the corresponding selected posterior
endpoint.  In particular its trial functional satisfies $G_k/k\le U$.
The hidden list is omitted when $h=0$, and the old descendant list is
omitted when $d=1$.
\end{lemma}

\begin{proof}
First fix all packet counts and signs, so the endpoint uses finitely
many distinct sites and its loss $X$ is bounded, say by $C$.
For each fixed $B$, the empirical tree evaluation is a continuous
function of finitely many retained message values and of the old
exponents.  Continuity follows also at messages equal to zero or
one, because every factor $1-wf_1f_2f_3$ is at least
$1-w=e^{-\beta}>0$.  The old exponents converge to positive limits
since $t_j\ge m_0$.
Thus weak convergence of the natural-order arrays gives convergence
of the expected empirical evaluation.  Lemma~\ref{ct:empirical}
is uniform before the limit, so taking $B\to\infty$ afterward
recovers exactly the subsequential expected endpoint value.

We now evaluate that same limit using the representation and its
independent cut partition.  Below the sample level, replace
empirical evaluations by the true conditional transforms
$T_{t_2},\ldots,T_{t_d}$.  The error tends to zero by the same
uniform estimate, even after propagation through the top empirical
average.  Write $Y_{\alpha,l}\in[-C,C]$ for the resulting value at
sample position $(\alpha,l)\in\mathbb N^h\times\mathbb N$.
Condition on the full cut partition and on all represented variables
through the hidden levels, including their root variables.  The
remaining sample-level variables are independent across sample
positions, and their laws are identical within each hidden leaf.
If $\alpha(u)$ is the leaf containing natural sample $u$, then
the conditional mean of the first-$B$ natural-order average is
\[
  \frac1B\sum_{u=1}^B
          \E_{\rm sample}e^{mY_{\alpha(u),1}}.
\]
Its conditional variance is $O_C(B^{-1})$.  The empirical
frequencies of the hidden leaves converge to their masses
$(v_\alpha)$; convergence is in $\ell^1$, since these are
probability masses whose pointwise limits have sum one.  Hence,
in $L^1$,
\begin{equation}\label{ct:natural-average}
 \frac1B\sum_{u=1}^B e^{mY_{\alpha(u),l(u)}}
   \longrightarrow
       \sum_\alpha v_\alpha\E_{\rm sample}e^{mY_{\alpha,1}}.
\end{equation}
Boundedness makes passage to the divided logarithm legitimate.
For $h=0$ the sum consists of one term of weight one.

Set $V_\alpha=T_m^{\rm sample}Y_{\alpha,1}$.  The represented marks
and their root variables were chosen independently of the full partition,
so the RPC weights retain their original law conditional on those root
variables and are independent of the attached marks.  If $h=0$, the
divided logarithm in \eqref{ct:natural-average} is simply
$V_{\varnothing}$, where $\varnothing$ is the unique empty label.
If $h\ge1$, the hidden parameters satisfy
$0<x_1<\cdots<x_h<1$, exactly the strict cascade hypothesis of
Lemma~\ref{ctr:rpc}.  Applying its logarithmic identity conditionally on
the root data to marks $mV_\alpha$ shows that
the expectation of the divided logarithm on the right of
\eqref{ct:natural-average} is the nested transform with hidden
exponents $mx_1,\ldots,mx_h$, followed by the sample exponent $m$.
The scaling is the identity
\[
                            T_x(mV)=mT_{mx}V,
\]
applied successively at the hidden levels and followed by division
by $m$.  Below them remain the old exponents $t_2,\ldots,t_d$.
This proves precisely \eqref{ct:new-exponents}.  The hidden exponents,
when present, are increasing and lie strictly between zero and $m$.
Every old descendant exponent is at least $m$, and their old order
persists.  Thus these are admissible intermediate trial exponents,
including when some old exponents coincide or equal one.

Finally remove the conditioning on counts and signs.  Both the
prelimit and limit transforms are bounded in absolute value by
$k\log2+\beta\sum D_i$, with the fixed Poisson packet laws.
Their tails are uniformly integrable.  Truncating counts, applying
the proved finite-count assertion, and then removing the truncation
establishes the assertion for each expected endpoint.  Subtracting
the B endpoint from the A endpoint and using
\eqref{ct:selected-value} yields $G_k/k\le U$.
\end{proof}

\subsection{Passing the pathwise structural records}\label{ct:sub-records}

The endpoint calculation has supplied a trial with the required value
bound.  It remains to preserve the old records and to construct the new
one.  We will pass empirical products of site tests through the weak
limit; the following Hilbert-space fact then converts those pairwise
bounds into statements about almost every individual global path.

\begin{lemma}[Closed support bounds]\label{ct:support}
Let $V,V'$ be independent random variables with the same law in a
separable Hilbert space, and let $\delta,\varepsilon\ge0$.
\begin{enumerate}
\item If $|\langle V,V'\rangle|\le\delta^2$ almost surely, then
      $\|V\|\le\delta$ almost surely.
\item If $\langle V,V'\rangle$ almost surely belongs to a fixed closed
      interval of length $\varepsilon$, then the support of their
      common law has diameter at most $\sqrt{2\varepsilon}$.  Consequently
      $\|V-\E V\|\le\sqrt{2\varepsilon}$ almost surely whenever
      the expectation exists.
\end{enumerate}
If $V,V'$ are conditionally independent with the same conditional law
given a standard Borel random element, the corresponding conclusions
hold for almost every conditional law in which the stated pair bound holds.
\end{lemma}

\begin{proof}
A closed set of full product measure contains the product of the
two supports: otherwise a pair of support points has product
neighborhoods of positive measure disjoint from that set.
The first assertion follows by putting the same support point in
both coordinates.  In the second, for support points $v,v'$, both
squared norms and their inner product belong to the same interval,
so
\[
 \|v-v'\|^2=\|v\|^2+\|v'\|^2-2\langle v,v'\rangle
                              \le2\varepsilon.
\]
Jensen's inequality gives the assertion about the barycenter.
The same reasoning applies after disintegration; separability
ensures that the relevant supports have full measure.
\end{proof}

\begin{lemma}[Preservation of the old records]\label{ct:old-records}
Every old processed block satisfies Equation~\eqref{ctr:record}
in the represented hierarchy, with the same test family and
tolerance and with its boundaries shifted by $h$.
\end{lemma}

\begin{proof}
Let $(b,c]$ be an old block.  Since the input hierarchy is active,
$b\ge1$.  In the prelimit natural array we may therefore choose
two paths in the same first-level sample, agreeing through $b$
and diverging immediately afterward.  At site $i$, form the
product of their two values of
\[
                              \phi(S_c)-S_b(\phi),
                         \qquad \phi\in\Phi_j.
\]
Condition on the complete globals of the two paths.  The sites
are independent and identically distributed, and the conditional
mean of this product is the $L^2(\dd u_{0:b})$ inner product of
the two error functions
\[
    E_z(u_{0:b})=
          \int\phi(S_c)\,\dd u_{b+1:c}-S_b(\phi).
\]
The fresh site variables after the branching point are independent
on the two paths; their shared site variables are exactly those
through $b$.  Equation~\eqref{ctr:record}, which survived the
posterior and Gaussian changes of law, bounds both error norms
by $\delta$.  Cauchy--Schwarz therefore bounds the absolute
conditional mean by $\delta^2$.

Let $M_I$ be the empirical mean over the first $I$ sites.  The
summands belong to $[-1,1]$, so
\begin{equation}\label{ct:old-empirical-bound}
     \E\operatorname{dist}(M_I,[-\delta^2,\delta^2])^2\le I^{-1}.
\end{equation}
For fixed $I$, the integrand on its left is a bounded continuous function of finitely
many retained type values: $\phi$ is continuous, and evaluation
of a probability measure against $\phi$ is continuous in the
type-space topology.  The inequality thus passes to the limit.
On $I=2^j$, Markov's inequality and Borel--Cantelli imply that
the displayed distance tends to zero almost surely.  This holds
simultaneously for all countably many path choices and all the
finitely many tests.  It consequently remains true for the
partition-selected paths after reindexing.

In the product-tree representation, the conditional strong law
for the sites identifies the limit of $M_I$ with the inner
product of the analogous error vectors, now in
$L^2(\dd u_{0:b+h})$.  Conditional on the common globals through
the shifted block start, those two vectors are independent and
identically distributed as their continuation globals vary.
Thus their inner product has absolute value at most $\delta^2$
almost surely.  Apply the first assertion of
Lemma~\ref{ct:support} to this conditional law.  Each error
vector has norm at most $\delta$ for almost every global
continuation.  Fubini's theorem gives precisely
Equation~\eqref{ctr:record} for almost every complete global
path.  The weak-limit step used a closed finite-column condition;
the pathwise conclusion then followed from the conditional support bound.
\end{proof}

Thus every previously processed block retains its record.  For the new
block, the narrow support band will provide the corresponding pairwise
bound without a preexisting record.

\begin{lemma}[A new processed block]\label{ct:new-record}
Let $c=h+1$ be the sample level.  Let $b=1$ if the parent cut is
present, and let $b=0$ otherwise.  There is an ancestor-measurable
probability-valued type $S_b$ such that $(b,c]$ satisfies
Equation~\eqref{ctr:record} for the specified new test family
and tolerance $\delta$.  Its exponents lie in $[m-\eta,m]$.
If the parent cut is present, its exponent belongs to
$(0,(1-\eta)m]$.
\end{lemma}

\begin{proof}
For two distinct natural first-level samples, average over the
first $I$ sites the product of their $\phi_j(S_1)$ values.
Conditional on their globals and Gaussian fields, the site's
root variable is shared between the two samples, and its
first-level variables are independent.  The conditional mean
is exactly $Q_j$ from \eqref{ct:features}.  Therefore
\begin{equation}\label{ct:new-empirical-bound}
       \E\left|\frac1I\sum_{i=1}^I
              \phi_j(S_1^{1,i})\phi_j(S_1^{2,i})-Q_j(1,2)
          \right|^2\le I^{-1}.
\end{equation}
For each fixed $I$ this inequality passes by continuity to the
natural-order limit array.  Along dyadic $I$, it gives almost-sure
convergence of the empirical average to $Q_j$, simultaneously
for all pairs and tests.

After reindexing, take two sample positions agreeing through $b$
and diverging immediately below $b$.  If the parent is present,
they are in the same parent cluster; otherwise $\Prob(Q<q)=0$.
If the child is present, they lie in distinct child clusters;
otherwise $\Prob(Q<r)=1$.  Consequently their natural pair
overlap satisfies $q\le Q<r$ almost surely in all cases.
By \eqref{ct:quantile-band} and
Lemma~\ref{ct:ordering}, the corresponding $Q_j$ belongs to
a fixed closed interval $I_j$ of length at most $J_f\gamma$.
The empirical characterization just proved shows that this
interval condition is a property of the reindexed decorations
alone; we do not need a separate representation of the kernels
$Q_j$.

In the product-tree representation, the site empirical limit is
the inner product in $L^2(\dd u_{0:b})$ of the two vectors
\begin{equation}\label{ct:new-vectors}
          A_j(u_{0:b})=
                     \int\phi_j(S_c)\,\dd u_{b+1:c}.
\end{equation}
Conditional on the ancestor globals through $b$, these are
independent copies as their continuation globals vary.
Their inner product belongs to $I_j$ almost surely.
The second assertion of Lemma~\ref{ct:support} shows that
the conditional support of $A_j$ has diameter at most
$\sqrt{2J_f\gamma}$.

Define, for each ancestor global and site history,
\begin{equation}\label{ct:new-type}
 S_b(z_{0:b},u_{0:b})
       =\operatorname{Law}_{z_{b+1:c},u_{b+1:c}}(S_c).
\end{equation}
All continuation variables in this formula have their current
independent uniform laws.  This is a measurable probability
kernel on the type space of $S_c$, hence has exactly the required
next type space.  By Fubini's theorem,
$S_b(\phi_j)$ is the conditional continuation-global barycenter
of the vector $A_j$ in \eqref{ct:new-vectors}.  The support
diameter bound gives
\[
  \int\left|
           \int\phi_j(S_c)\,\dd u_{b+1:c}-S_b(\phi_j)
      \right|^2\dd u_{0:b}\le 2J_f\gamma
\]
for almost every complete global path.  Choose
$0<\gamma\le\delta^2/(2J_f)$ to obtain the new record.

It remains to check the exponent band.  The only possible
intermediate level between $b$ and $c$ is the child cut, whose
exponent is $mx_c$.  By \eqref{ct:quantile-band},
$x_c\ge1-\eta$, so $m-mx_c\le m\eta\le\eta$.
The end exponent is $m$.  If the parent is present, its exponent
satisfies $0<mx_p\le(1-\eta)m$, as claimed.

The type in \eqref{ct:new-type} is defined once, at this
transition.  In subsequent posterior changes it is retained as
the same measurable function; it is not redefined by averaging
under the changed global law.  The pathwise record just proved
then survives by absolute continuity, exactly as do the old
records.
\end{proof}

\begin{proof}[Completion of the proof of Lemma~\ref{ct:transition}]
The compactness and representation construction inserts at most
two levels.  Lemma~\ref{ct:endpoint} supplies the limiting expected
endpoint values and $G_k/k\le U$.
Lemma~\ref{ct:old-records} preserves every old record, while
their exponent-band inequalities survive ordinary limits.
Lemma~\ref{ct:new-record} supplies the new processed block.
If there is no parent cut it begins at the root and processing
is complete.  Otherwise the hierarchy is active, its first
exponent is $mx_p>0$, and this is at most $(1-\eta)m$.
All records hold for almost every complete global path.  These constructions were made for the fixed output
horizon $k$; repeating them for any other fixed horizon proves
the asserted quantifiers.  All nonroot shared globals remain in this
output.  The next section removes them by a separate approximation
argument.
\end{proof}
